\section{Anticipar lo que anticiparán las grandes empresas: espera, tentación y punto de inflexión}

La sección anterior trató el capital y la comercialización; esta examina la capacidad táctica más importante de Xiao Hong (肖弘): anticipar lo que anticiparán las grandes empresas. La oportunidad del SCRM para WeCom (企业微信) surgió de un juicio: Tencent (腾讯) iba a regular los complementos no oficiales de WeChat personal y, al mismo tiempo, iba a canalizar la demanda de cumplimiento normativo hacia el ecosistema de WeCom. No se sabía exactamente cuándo ocurriría, pero era posible prepararse con anticipación dentro de un costo asumible. Cuando se combatieron los complementos no oficiales, el producto ya preparado pudo atender de inmediato esa demanda.

\begin{figure}[H]
\centering
\includegraphics[width=0.96\textwidth]{predict-giants.png}
\caption{Anticipar lo que anticiparán las grandes empresas: la espera, la tentación y el anticonsenso al estilo de «La gran apuesta» forman en conjunto el punto de inflexión decisivo. Diagrama conceptual propio, elaborado a partir de la conversación entre 00:42:10 y 00:52:27.}
\end{figure}

\begin{knowledgebox}{Lectura de la figura: se sabe que ocurrirá, pero no cuándo}
La figura avanza de la observación a la anticipación, la espera, la tentación y el punto de inflexión. El método central de Xiao Hong es este: si juzgas que algo ocurrirá con alta probabilidad, pero no conoces el momento, prepárate con anticipación dentro de un costo asumible. La dificultad del emprendimiento está en que durante la espera no hay retroalimentación positiva, y además aparecen proyectos a la medida, pedidos de corto plazo y titubeos en el equipo.
\end{knowledgebox}

\subsection{La espera al estilo de «La gran apuesta»: por qué es lo más difícil}

Esta subsección explica lo doloroso de la espera. Después del lanzamiento del producto SCRM, los complementos no oficiales todavía no se habían regulado, el producto no tenía usuarios y el equipo tendía a dudar de la estrategia. Al mismo tiempo surgían tentaciones: desarrollos a la medida para grandes clientes, ingresos de corto plazo y otras direcciones. Xiao Hong recurre a la analogía de «La gran apuesta»: acertar en el juicio no equivale a ganar dinero de inmediato; durante la espera hay que soportar presión de efectivo, de equipo y de confianza. Finalmente, cuando se bloquearon los complementos no oficiales, esa misma noche escribieron un artículo, lo difundieron con publicidad pagada y se posicionaron rápidamente en la mente de los usuarios como la alternativa.

\begin{warningbox}{El riesgo del periodo de espera: no es equivocarse en el juicio, sino no resistir hasta que ocurra}
Los emprendedores con frecuencia alcanzan a ver una dirección, pero no el momento. El juicio sobre la dirección debe corresponder con el ciclo de efectivo, la paciencia del equipo y el grado de terminación del producto. Llegar demasiado pronto desgasta; llegar demasiado tarde hace perder la ventana.
\end{warningbox}

\subsection{Resistir la retroalimentación positiva ajena a la estrategia}

Esta subsección añade un detalle. Durante la espera, los desarrollos a la medida para grandes clientes pueden dar retroalimentación positiva al equipo; el equipo de ventas se entusiasma y el efectivo de corto plazo también resulta atractivo. Pero si los proyectos a la medida se alejan del producto SaaS estándar, absorben recursos de desarrollo y afectan la iteración de la línea principal. Lo que hizo Xiao Hong fue rechazar la retroalimentación positiva ajena a la estrategia y reservar los recursos para el producto que «pudiera atender la demanda cuando ocurriera el evento». Este juicio es más concreto que el «enfoque» habitual: no se trata de rechazar todos los ingresos, sino de rechazar los ingresos que destruirían la espera estratégica.

\begin{importantbox}{Experiencia práctica: no toda retroalimentación positiva debe aceptarse}
Si una retroalimentación proviene de una dirección no estratégica, puede ser más peligrosa que una retroalimentación negativa. La retroalimentación negativa te pone en alerta; la retroalimentación positiva ajena a la estrategia lleva al equipo a creer erróneamente que encontró un nuevo camino y termina desviándolo de la ventana que en realidad estaba esperando.
\end{importantbox}

\subsection{Resumen de la sección}

Anticipar lo que anticiparán las grandes empresas fue el método central del primer emprendimiento de Xiao Hong: entender cómo regulará una plataforma su ecosistema, preparar con anticipación una alternativa que cumpla con la normativa y resistir las tentaciones de corto plazo durante la espera. Este método se trasladó después a la IA: anticipar cuándo aparecerán las capacidades de los modelos y colocar de antemano ahí el contenedor de la aplicación.
